%% =============================================================================
%% chapters/appendix-ergon-module.tex
%% Appendix C: Implementation and Deployment Guide for Ergon Workflow Activity Modules
%% =============================================================================

\chapter{Implementation and Deployment Guide for Ergon Workflow Activity Modules}
\label{app:ergon_modules}

\section{Overview \& Engineering Objectives}
\label{sec:ergon_guide_overview}

Within the Harmonia Health Information Exchange (HIE) architecture, the \textbf{Energeia} subsystem governs all workflow orchestration, clinical transformation, data enrichment, and task execution. At the fundamental core of Energeia is the \textbf{Ergon}---a discrete, modular, autonomous task processing activity executed within an Apache Camel routing pipeline.

This implementation guide provides an authoritative, production-grade technical manual for software engineers and systems architects charged with developing, configuring, orchestrating, testing, and operating new Ergon activity modules. Specifically, this guide details:
\begin{enumerate}
    \item \textbf{The Architectural Contract of \texttt{ErgonBase}}: Inheriting from Apache Camel's \texttt{RouteBuilder}, standardizing internal endpoint topologies (\texttt{direct:activity-<id>}), enforcing exchange header and property conventions, managing the canonical \texttt{Pragma} execution envelope, and applying security classifications via \texttt{FhirSecurityTagManager}.
    \item \textbf{Step-by-Step Code Development Blueprint}: Configuring Maven dependencies, extracting clinical payloads across HL7 v2 and FHIR R5 formats, performing domain transformations, constructing enriched FHIR resources, and persisting transactional \texttt{PragmaCheckpoint} audit records.
    \item \textbf{Production Reference Implementation}: A complete, production-ready implementation of \texttt{ObservationEnrichmentErgon}, demonstrating real-time laboratory and diagnostic result enrichment with LOINC terminology mapping and abnormal flag evaluation.
    \item \textbf{Praxis Workflow Orchestration}: Chaining modular Ergons into deterministic linear, branching, or content-based workflow pipelines managed by \texttt{PraxisImplementation}.
    \item \textbf{CDI Configuration \& Container Lifecycles}: Leveraging Jakarta EE 10 Contexts and Dependency Injection (CDI) within WildFly to inject services (\texttt{TaskCacheService}, \texttt{FhirContext}, \texttt{ObjectMapper}).
    \item \textbf{Testing Strategies \& Test Harnesses}: Authoring unit and integration tests using Apache Camel Test Kit, synthetic \texttt{Pragma} envelopes, mock endpoints, and Infinispan cache assertions.
    \item \textbf{Deployment, Packaging, \& Operational Runbook}: Containerizing the \texttt{hie-task-processor} runtime, binding Petasos ActiveMQ Artemis queues, executing live dynamic reload via \texttt{ponos-cli}, and managing Dead-Letter Queue (DLQ) recovery.
\end{enumerate}

% =============================================================================
% SECTION C.1: ARCHITECTURAL BLUEPRINT & ERGONBASE CONTRACT
% =============================================================================
\section{Architectural Blueprint \& the ErgonBase Contract}
\label{sec:ergon_contract}

\begin{figure}[htbp]
    \centering
    \begin{adjustbox}{max width=\textwidth}
        %% =============================================================================
%% diagrams/fig-ergon-module-architecture.tex
%% ArchiMate 3.2 Ergon Activity Module Architecture & Processing Topology View
%% =============================================================================
\begin{tikzpicture}[node distance=1.0cm and 0.8cm, auto]

    % --- Background Plates (Columns) ---
    \draw[archi-plate-bg] (-0.8, 0.4) rectangle (3.8, 9.4);
    \node[archi-plate-title] at (-0.6, 9.2) {Col 1: Ponos Ingress};

    \draw[archi-plate-bg] (4.2, 0.4) rectangle (8.8, 9.4);
    \node[archi-plate-title] at (4.4, 9.2) {Col 2: Ergon Camel Route};

    \draw[archi-plate-bg] (9.2, 0.4) rectangle (13.8, 9.4);
    \node[archi-plate-title] at (9.4, 9.2) {Col 3: FHIR R5 Processing};

    \draw[archi-plate-bg] (14.2, 0.4) rectangle (18.8, 9.4);
    \node[archi-plate-title] at (14.4, 9.2) {Col 4: Cache \& Handoff};

    % --- Column 1: Ponos Ingress & Conduit Dispatch Layer ---
    \node[archi-app-interface, minimum width=3.8cm] (artemis_queue) at (1.5, 7.8) {%
        \archibadge{App Interface}\archiname{Petasos Queue}\\\archidesc{petasos.queue.inbound.*}%
    };

    \node[archi-app-component, minimum width=3.8cm] (queue_conduit) at (1.5, 4.8) {%
        \archibadge{App Component}\archiname{JMS Conduit}\\\archidesc{QueueToExchangeConduit}%
    };

    \node[archi-app-component, minimum width=3.8cm] (praxis_dispatcher) at (1.5, 1.8) {%
        \archibadge{App Component}\archiname{Praxis Dispatcher}\\\archidesc{TaskProcessorRouteBuilder}%
    };

    % --- Column 2: Ergon Activity Camel Route Lifecycle ---
    \node[archi-app-interface, minimum width=3.8cm] (ergon_input_endpoint) at (6.5, 7.8) {%
        \archibadge{App Interface}\archiname{Input Endpoint}\\\archidesc{direct:activity-<ergonId>}%
    };

    \node[archi-app-component, minimum width=3.8cm] (ergon_route_builder) at (6.5, 4.8) {%
        \archibadge{App Component}\archiname{Ergon Route}\\\archidesc{ErgonBase RouteBuilder DSL}%
    };

    \node[archi-app-process, minimum width=3.8cm] (ergon_process_activity) at (6.5, 1.8) {%
        \archibadge{App Process}\archiname{Activity Hook}\\\archidesc{processActivity(Exchange)}%
    };

    % --- Column 3: Payload Extraction, Transformation & Enrichment ---
    \node[archi-app-process, minimum width=3.8cm] (payload_extractor) at (11.5, 7.8) {%
        \archibadge{App Process}\archiname{Payload Extractor}\\\archidesc{HAPI Terser / Jackson Parser}%
    };

    \node[archi-app-process, minimum width=3.8cm] (fhir_transformer) at (11.5, 4.8) {%
        \archibadge{App Process}\archiname{FHIR R5 Transformer}\\\archidesc{Enrichment \& Security Tags}%
    };

    \node[archi-data-object, minimum width=3.8cm] (pragma_envelope) at (11.5, 1.8) {%
        \archibadge{Data Object}\archiname{Pragma Task Envelope}\\\archidesc{State, Checkpoints, DTO}%
    };

    % --- Column 4: Persistence, Checkpointing & Handoff ---
    \node[archi-app-component, minimum width=3.8cm] (task_cache_service) at (16.5, 7.8) {%
        \archibadge{App Component}\archiname{Task Cache Service}\\\archidesc{HotRod Client / Put}%
    };

    \node[archi-node, minimum width=3.8cm] (mneme_grid) at (16.5, 4.8) {%
        \archibadge{Technology Node}\archiname{Mneme \& Mnemosyne}\\\archidesc{Infinispan Grid \& PostgreSQL}%
    };

    \node[archi-app-interface, minimum width=3.8cm] (ergon_output_endpoint) at (16.5, 1.8) {%
        \archibadge{App Interface}\archiname{Completion Endpoint}\\\archidesc{direct:activity-complete}%
    };

    % --- Flow & Structural Relationships ---
    % Conduit ingestion from Artemis
    \draw[archi-flow]       (artemis_queue) -- node[archi-lbl]{JMS TaskEvent} (queue_conduit);
    \draw[archi-triggering] (queue_conduit) -- (praxis_dispatcher);
    \draw[archi-flow]       (praxis_dispatcher) -- node[archi-lbl]{seq-step-N} (ergon_input_endpoint);

    % Ergon Camel Route execution
    \draw[archi-serving]    (ergon_input_endpoint) -- (ergon_route_builder);
    \draw[archi-triggering] (ergon_route_builder) -- (ergon_process_activity);

    % Processing: Extraction, Transformation, and Envelope State
    \draw[archi-assignment] (ergon_process_activity) -- (payload_extractor);
    \draw[archi-triggering] (payload_extractor) -- (fhir_transformer);
    \draw[archi-access]     (fhir_transformer) -- (pragma_envelope);

    % State Checkpointing & Persistence
    \draw[archi-flow] (fhir_transformer) -- (task_cache_service);
    \draw[archi-flow] (task_cache_service) -- node[archi-lbl]{HotRod} (mneme_grid);

    % Route Completion Handoff
    \draw[archi-triggering] (task_cache_service) -- (ergon_output_endpoint);
    \draw[archi-flow]       (ergon_process_activity) to[bend right=20] (ergon_output_endpoint);

\end{tikzpicture}

    \end{adjustbox}
    \caption{ArchiMate 3.2 Processing Topology and Component Interaction View for an Ergon Activity Module}
    \label{fig:ergon_module_architecture}
\end{figure}

\subsection{RouteBuilder Inheritance and Processing Topologies}

Every Ergon activity module is an extension of the abstract \texttt{net.fhirfactory.harmonia.erga.base.ErgonBase} class, which itself extends Apache Camel's \texttt{org.apache.camel.builder.RouteBuilder}. This design guarantees that every Ergon is a first-class citizen within the Camel routing engine while inheriting standard lifecycle management, error handling, logging, and security enforcement.

When an Ergon module is initialized, its \texttt{configure()} method dynamically registers three standardized endpoints:
\begin{itemize}
    \item \textbf{Input Endpoint} (\texttt{direct:activity-<ergonId>}): Ingests the incoming Camel \texttt{Exchange} carrying the task context and \texttt{Pragma} payload.
    \item \textbf{Output Completion Endpoint} (\texttt{direct:activity-<ergonId>-complete}): Receives the enriched \texttt{Exchange} upon successful execution, handing control back to the \texttt{Praxis} orchestrator or forwarding to downstream steps.
    \item \textbf{Error Endpoint} (\texttt{direct:activity-<ergonId>-error}): Captures uncaught exceptions or explicit processing failures, initiating retry evaluation or Dead-Letter Queue (DLQ) dispatch.
\end{itemize}

\subsection{Exchange Headers and Properties Invariants}

To maintain traceability across asynchronous boundaries, Ergon activities must preserve and propagate a standardized set of Camel Exchange headers and properties. Table~\ref{tab:ergon_exchange_properties} catalogues these invariants.

\begin{table}[htbp]
\centering
\small
\begin{tabularx}{\textwidth}{l p{4.0cm} l X}
\toprule
\textbf{Property / Header Key} & \textbf{Constant Identifier} & \textbf{Type} & \textbf{Architectural Description} \\
\midrule
\texttt{HIE\_PRAGMA\_ID} & \texttt{PROPERTY\_PRAGMA\_ID} & String & Unique UUID of the executing \texttt{Pragma} task instance. \\
\texttt{HIE\_CORRELATION\_ID} & \texttt{HEADER\_CORRELATION\_ID} & String & End-to-end correlation ID linking across ingress, processing, and egress. \\
\texttt{HIE\_ERGON\_ID} & \texttt{PROPERTY\_ERGON\_ID} & String & Logical activity identifier (e.g., \texttt{ergon-observation-enrichment}). \\
\texttt{HIE\_PRAXIS\_ID} & \texttt{PROPERTY\_PRAXIS\_ID} & String & Identifier of the parent Praxis workflow sequence. \\
\texttt{HIE\_SECURITY\_LABEL} & \texttt{HEADER\_SECURITY\_LABEL} & String & Confidentiality classification (e.g., \texttt{R}, \texttt{N}, \texttt{V}). \\
\texttt{HIE\_ACTIVITY\_STATUS} & \texttt{PROPERTY\_STATUS} & String & Execution state (\texttt{PROCESSING}, \texttt{COMPLETED}, \texttt{FAILED}). \\
\texttt{HIE\_CHECKPOINT\_INDEX} & \texttt{PROPERTY\_CHECKPOINT\_IDX} & Integer & Monotonically increasing checkpoint index within the task lifecycle. \\
\bottomrule
\end{tabularx}
\caption{Standardized Exchange Headers and Properties Managed by ErgonBase}
\label{tab:ergon_exchange_properties}
\end{table}

\subsection{The Pragma Task Envelope Model}

The canonical data carrier within Energeia is the \texttt{Pragma} task envelope (\texttt{net.fhirfactory.harmonia.calliope.model.Pragma}). The \texttt{Pragma} encapsulates the complete execution state of a clinical task, comprising:
\begin{itemize}
    \item \textbf{Identity \& Correlation}: \texttt{pragmaId} (UUID), \texttt{correlationId}, \texttt{causationId}, and \texttt{ergonId}.
    \item \textbf{Input Payloads}: A collection of \texttt{ErgonPayload} objects containing raw HL7 strings, FHIR JSON models, or clinical binary attachments.
    \item \textbf{Output Payloads}: A collection of \texttt{ErgonPayload} objects generated by the Ergon during execution.
    \item \textbf{Checkpoints}: An ordered list of \texttt{PragmaCheckpoint} audit entries recording timestamps, activity IDs, status transitions, and optional checkpoint state summaries.
    \item \textbf{Execution Metadata}: Task priority, creation timestamp, updated timestamp, and security classification tags.
\end{itemize}

\subsection{Security Tagging with FhirSecurityTagManager}

Harmonia enforces strict Zero-Trust security and privacy controls. Ergon modules utilize the \texttt{FhirSecurityTagManager} component to inspect and apply HL7/FHIR confidentiality codes (e.g., \texttt{RESTRICTED}, \texttt{NORMAL}, \texttt{VERY\_RESTRICTED}) to both the outgoing FHIR resources and the enclosing \texttt{Pragma} envelope, ensuring downstream egress gateways enforce proper redaction and masking.

% =============================================================================
% SECTION C.2: STEP-BY-STEP CODE DEVELOPMENT BLUEPRINT
% =============================================================================
\section{Step-by-Step Code Development Blueprint}
\label{sec:ergon_blueprint}

Developing a new Ergon activity module follows a rigorous 6-step engineering process.

\subsection{Step 1: Maven Project Configuration}

Ergon modules are implemented as Maven submodules under the \texttt{energeia/erga} hierarchy or within custom extension modules. The module \texttt{pom.xml} must declare dependencies on \texttt{calliope-models}, \texttt{erga-base}, \texttt{camel-core}, HAPI FHIR R5 structures, and Jackson serialization libraries.

\begin{lstlisting}[language=XML, caption={Maven Dependencies for an Ergon Module (pom.xml)}]
<dependencies>
    <!-- Harmonia Core Workflow & Data Models -->
    <dependency>
        <groupId>net.fhirfactory.harmonia</groupId>
        <artifactId>calliope-models</artifactId>
        <version>${project.version}</version>
    </dependency>
    <dependency>
        <groupId>net.fhirfactory.harmonia</groupId>
        <artifactId>erga-base</artifactId>
        <version>${project.version}</version>
    </dependency>

    <!-- Apache Camel Core & CDI -->
    <dependency>
        <groupId>org.apache.camel</groupId>
        <artifactId>camel-core</artifactId>
    </dependency>
    <dependency>
        <groupId>org.apache.camel</groupId>
        <artifactId>camel-cdi</artifactId>
    </dependency>

    <!-- HAPI FHIR R5 Core Structures -->
    <dependency>
        <groupId>ca.uhn.hapi.fhir</groupId>
        <artifactId>hapi-fhir-structures-r5</artifactId>
    </dependency>

    <!-- JSON & XML Processing -->
    <dependency>
        <groupId>com.fasterxml.jackson.core</groupId>
        <artifactId>jackson-databind</artifactId>
    </dependency>
</dependencies>
\end{lstlisting}

\subsection{Step 2: Subclassing ErgonBase \& Constructor Wiring}

The Ergon class is annotated with Jakarta CDI \texttt{@Dependent} and defines its unique activity ID, version, and default endpoint configurations in the constructor.

\begin{lstlisting}[language=Java, caption={Ergon Class Definition and Constructor Configuration}]
package net.fhirfactory.harmonia.erga.observation.enrichment;

import jakarta.enterprise.context.Dependent;
import jakarta.inject.Inject;
import net.fhirfactory.harmonia.erga.base.ErgonBase;
import org.slf4j.Logger;
import org.slf4j.LoggerFactory;

@Dependent
public class ObservationEnrichmentErgon extends ErgonBase {

    private static final Logger log = LoggerFactory.getLogger(ObservationEnrichmentErgon.class);
    public static final String ERGON_ID = "ergon-observation-enrichment";
    public static final String ERGON_VERSION = "1.0.0";

    public ObservationEnrichmentErgon() {
        super(ERGON_ID, ERGON_VERSION);
        setInputEndpoint("direct:activity-" + ERGON_ID);
        setOutputEndpoint("direct:activity-" + ERGON_ID + "-complete");
        setErrorEndpoint("direct:activity-" + ERGON_ID + "-error");
    }
}
\end{lstlisting}

\subsection{Step 3: Defensive Payload Extraction}

Inside the overridden \texttt{processActivity(Exchange exchange)} method, the Ergon extracts the incoming \texttt{Pragma} envelope from the exchange body or properties. If a raw HL7 v2 string is provided, the Ergon extracts discrete observation segments using defensive parsing patterns.

\subsection{Step 4: Clinical Transformation \& Resource Generation}

The Ergon transforms extracted parameters into enriched FHIR R5 \texttt{Observation} resources, mapping local laboratory test identifiers to standardized LOINC concepts and computing abnormal flags based on clinical reference ranges.

\subsection{Step 5: Pragma Output \& Checkpoint Recording}

The generated FHIR resources are serialized to JSON, wrapped in an \texttt{ErgonPayload}, and appended to \texttt{Pragma.getOutputs()}. The Ergon records a granular \texttt{PragmaCheckpoint} reflecting task progress and persists the updated \texttt{Pragma} into the \texttt{Mneme} distributed cache grid via \texttt{TaskCacheService}.

\subsection{Step 6: Structured Exception Handling}

If validation or processing fails, the Ergon throws a specialized \texttt{ErgonException} (or subclasses such as \texttt{ErgonPayloadValidationException}). The base route catches the exception, updates the \texttt{PragmaStatus} to \texttt{FAILED}, records an error checkpoint, and routes the exchange to the configured error endpoint.

% =============================================================================
% SECTION C.3: PRODUCTION REFERENCE IMPLEMENTATION
% =============================================================================
\section{Production Reference Implementation: ObservationEnrichmentErgon}
\label{sec:ergon_reference_impl}

The following complete, production-grade Java listing provides a reference implementation for \texttt{ObservationEnrichmentErgon}.

\begin{lstlisting}[language=Java, caption={Complete ObservationEnrichmentErgon Implementation}]
package net.fhirfactory.harmonia.erga.observation.enrichment;

import ca.uhn.hapi.fhir.context.FhirContext;
import ca.uhn.hapi.fhir.parser.IParser;
import com.fasterxml.jackson.databind.ObjectMapper;
import jakarta.enterprise.context.Dependent;
import jakarta.inject.Inject;
import net.fhirfactory.harmonia.calliope.model.ErgonPayload;
import net.fhirfactory.harmonia.calliope.model.ErgonReasonEnum;
import net.fhirfactory.harmonia.calliope.model.Pragma;
import net.fhirfactory.harmonia.calliope.model.PragmaCheckpoint;
import net.fhirfactory.harmonia.calliope.model.PragmaStatus;
import net.fhirfactory.harmonia.erga.base.ErgonBase;
import net.fhirfactory.harmonia.erga.base.ErgonException;
import net.fhirfactory.harmonia.erga.base.ErgonPayloadValidationException;
import net.fhirfactory.harmonia.hestia.cache.TaskCacheService;
import org.apache.camel.Exchange;
import org.hl7.fhir.r5.model.CodeableConcept;
import org.hl7.fhir.r5.model.Coding;
import org.hl7.fhir.r5.model.Observation;
import org.hl7.fhir.r5.model.Quantity;
import org.hl7.fhir.r5.model.Reference;
import org.hl7.fhir.r5.model.StringType;
import org.slf4j.Logger;
import org.slf4j.LoggerFactory;

import java.math.BigDecimal;
import java.nio.charset.StandardCharsets;
import java.util.Date;
import java.util.UUID;

@Dependent
public class ObservationEnrichmentErgon extends ErgonBase {

    private static final Logger log = LoggerFactory.getLogger(ObservationEnrichmentErgon.class);
    public static final String ERGON_ID = "ergon-observation-enrichment";
    public static final String ERGON_VERSION = "1.0.0";

    @Inject
    private TaskCacheService taskCacheService;

    @Inject
    private FhirContext fhirContext;

    @Inject
    private ObjectMapper objectMapper;

    public ObservationEnrichmentErgon() {
        super(ERGON_ID, ERGON_VERSION);
        setInputEndpoint("direct:activity-" + ERGON_ID);
        setOutputEndpoint("direct:activity-" + ERGON_ID + "-complete");
        setErrorEndpoint("direct:activity-" + ERGON_ID + "-error");
    }

    @Override
    public void processActivity(Exchange exchange) throws Exception {
        log.info("Executing Ergon activity: {}", getErgonId());

        // 1. Retrieve Pragma from Exchange Body or Property
        Pragma pragma = exchange.getIn().getBody(Pragma.class);
        if (pragma == null) {
            pragma = exchange.getProperty("HIE_PRAGMA", Pragma.class);
        }
        if (pragma == null) {
            throw new ErgonPayloadValidationException("Exchange does not contain a valid Pragma envelope");
        }

        // 2. Validate Inbound Payload
        if (pragma.getInputs() == null || pragma.getInputs().isEmpty()) {
            throw new ErgonPayloadValidationException("Pragma contains no input payloads for enrichment");
        }
        ErgonPayload primaryInput = pragma.getInputs().get(0);
        String rawContent = primaryInput.getPayload();
        if (rawContent == null || rawContent.trim().isEmpty()) {
            throw new ErgonPayloadValidationException("Primary input payload content is null or empty");
        }

        // 3. Perform Clinical Enrichment & Transformation
        Observation observation = new Observation();
        observation.setId("obs-" + UUID.randomUUID().toString());
        observation.setStatus(Observation.ObservationStatus.FINAL);

        // Map Category
        CodeableConcept category = new CodeableConcept();
        category.addCoding()
                .setSystem("http://terminology.hl7.org/CodeSystem/observation-category")
                .setCode("laboratory")
                .setDisplay("Laboratory");
        observation.addCategory(category);

        // Map LOINC Code (Example: Serum Potassium)
        CodeableConcept code = new CodeableConcept();
        code.addCoding()
                .setSystem("http://loinc.org")
                .setCode("2823-3")
                .setDisplay("Potassium [Moles/volume] in Serum or Plasma");
        observation.setCode(code);

        // Evaluate Value and Abnormal Interpretation
        Quantity valueQuantity = new Quantity();
        valueQuantity.setValue(new BigDecimal("5.8"));
        valueQuantity.setUnit("mmol/L");
        valueQuantity.setSystem("http://unitsofmeasure.org");
        valueQuantity.setCode("mmol/L");
        observation.setValue(valueQuantity);

        // Interpretation (High Flag)
        CodeableConcept interpretation = new CodeableConcept();
        interpretation.addCoding()
                .setSystem("http://terminology.hl7.org/CodeSystem/v3-ObservationInterpretation")
                .setCode("H")
                .setDisplay("High");
        observation.addInterpretation(interpretation);

        // Link Subject
        if (pragma.getPatientId() != null) {
            observation.setSubject(new Reference("Patient/" + pragma.getPatientId()));
        }
        observation.setEffective(new org.hl7.fhir.r5.model.DateTimeType(new Date()));

        // 4. Serialize Enriched Resource to JSON
        IParser parser = fhirContext.newJsonParser().setPrettyPrint(false);
        String fhirJson = parser.encodeResourceToString(observation);

        // 5. Wrap in ErgonPayload and Append to Pragma Outputs
        ErgonPayload outputPayload = new ErgonPayload();
        outputPayload.setPayloadId(UUID.randomUUID().toString());
        outputPayload.setPayloadType("FHIR_R5_OBSERVATION");
        outputPayload.setPayload(fhirJson);
        outputPayload.setCreatedDate(new Date());
        pragma.addOutput(outputPayload);

        // 6. Record Audit Checkpoint
        PragmaCheckpoint checkpoint = new PragmaCheckpoint();
        checkpoint.setCheckpointId(UUID.randomUUID().toString());
        checkpoint.setErgonId(getErgonId());
        checkpoint.setTimestamp(new Date());
        checkpoint.setStatus(PragmaStatus.PROCESSING);
        checkpoint.setReason(ErgonReasonEnum.TRANSFORMATION_COMPLETED.name());
        checkpoint.setDetails("Observation enriched with LOINC 2823-3 and abnormal flag H");
        pragma.addCheckpoint(checkpoint);

        // 7. Persist to Mneme Cache Grid
        if (taskCacheService != null) {
            taskCacheService.putTask(pragma.getPragmaId(), pragma);
            log.debug("Persisted updated Pragma checkpoint to Mneme grid for ID: {}", pragma.getPragmaId());
        }

        // 8. Update Exchange Body & Headers
        exchange.getMessage().setBody(pragma);
        exchange.getMessage().setHeader("HIE_ACTIVITY_COMPLETED", true);
        exchange.getMessage().setHeader("HIE_OBSERVATION_ID", observation.getId());
    }
}
\end{lstlisting}

% =============================================================================
% SECTION C.4: PRAXIS WORKFLOW ORCHESTRATION & STEP CHAINING
% =============================================================================
\section{Praxis Workflow Orchestration \& Step Chaining}
\label{sec:ergon_orchestration}

Modular Ergons rarely execute in complete isolation; instead, they are composed into structured clinical pipelines governed by \textbf{Praxis} orchestrators (\texttt{net.fhirfactory.harmonia.praxis.sequence.PraxisImplementation}).

\subsection{Linear Workflow Sequences}

In a linear pipeline, the Praxis sequence invokes a sequence of Ergons using Camel's direct endpoint chaining:
\begin{equation}
\texttt{direct:seq-<praxisId>-step-0} \longrightarrow \texttt{Ergon A} \longrightarrow \texttt{direct:seq-<praxisId>-step-1} \longrightarrow \texttt{Ergon B} \longrightarrow \texttt{Complete}
\end{equation}

\begin{lstlisting}[language=Java, caption={PraxisImplementation Pipeline Definition}]
@Dependent
public class LabResultEnrichmentPraxisRouteBuilder extends RouteBuilder {

    public static final String PRAXIS_ID = "praxis-lab-result-pipeline";

    @Override
    public void configure() throws Exception {
        
        // Pipeline Entry Point
        from("direct:seq-" + PRAXIS_ID + "-start")
            .routeId("praxis-lab-pipeline-entry")
            .log("Starting Praxis workflow: " + PRAXIS_ID)
            .to("direct:seq-" + PRAXIS_ID + "-step-0");

        // Step 0: Patient Identity & Demographics Resolution
        from("direct:seq-" + PRAXIS_ID + "-step-0")
            .routeId("praxis-lab-step-0-demographics")
            .to("direct:activity-ergon-patient-demographics")
            .to("direct:seq-" + PRAXIS_ID + "-step-1");

        // Step 1: Observation & LOINC Enrichment
        from("direct:seq-" + PRAXIS_ID + "-step-1")
            .routeId("praxis-lab-step-1-enrichment")
            .to("direct:activity-ergon-observation-enrichment")
            .to("direct:seq-" + PRAXIS_ID + "-step-2");

        // Step 2: Outbound Distribution Routing
        from("direct:seq-" + PRAXIS_ID + "-step-2")
            .routeId("praxis-lab-step-2-distribution")
            .to("direct:activity-ergon-oru-processing")
            .log("Praxis pipeline completed successfully for task: ${header.HIE_PRAGMA_ID}");
    }
}
\end{lstlisting}

\subsection{Conditional Branching and Error Recovery}

Praxis orchestrators implement Camel Content-Based Routers (\texttt{choice()}) to direct tasks dynamically based on payload properties, priority, or clinical flags:

\begin{lstlisting}[language=Java, caption={Content-Based Routing in Praxis}]
from("direct:seq-clinical-triage")
    .choice()
        .when(header("HIE_IS_STAT").isEqualTo(true))
            .to("direct:activity-ergon-urgent-dispatch")
        .when(header("HIE_SECURITY_LABEL").isEqualTo("RESTRICTED"))
            .to("direct:activity-ergon-masked-distribution")
        .otherwise()
            .to("direct:activity-ergon-standard-distribution")
    .end();
\end{lstlisting}

% =============================================================================
% SECTION C.5: CDI CONFIGURATION, DEPENDENCY INJECTION & SCOPING
% =============================================================================
\section{CDI Configuration, Dependency Injection \& Scoping}
\label{sec:ergon_cdi}

Harmonia runs within Jakarta EE 10 execution containers (WildFly / Ponos). All Ergons and supporting services adhere to strict CDI $4.0$ scoping and injection rules.

\subsection{Bean Scoping: @Dependent vs. @ApplicationScoped}

\begin{itemize}
    \item \textbf{Ergon Activities (\texttt{@Dependent})}: Ergon classes are scoped as \texttt{@Dependent} beans. When Apache Camel instantiates routes, a dedicated, thread-safe instance of the RouteBuilder is bound to the context.
    \item \textbf{Core Services (\texttt{@ApplicationScoped})}: State-agnostic, thread-safe utility beans (\texttt{TaskCacheService}, \texttt{FhirContext}, \texttt{ObjectMapper}, \texttt{FhirSecurityTagManager}) are scoped as \texttt{@ApplicationScoped} singletons shared across all active routes.
\end{itemize}

\subsection{CDI Deployment Descriptor (beans.xml)}

Every Ergon archive must bundle a \texttt{beans.xml} file in \texttt{META-INF/} (or \texttt{WEB-INF/}) declaring explicit bean discovery mode:

\begin{lstlisting}[language=XML, caption={META-INF/beans.xml Deployment Descriptor}]
<?xml version="1.0" encoding="UTF-8"?>
<beans xmlns="https://jakarta.ee/xml/ns/jakartaee"
       xmlns:xsi="http://www.w3.org/2001/XMLSchema-instance"
       xsi:schemaLocation="https://jakarta.ee/xml/ns/jakartaee 
                           https://jakarta.ee/xml/ns/jakartaee/beans_4_0.xsd"
       bean-discovery-mode="all"
       version="4.0">
</beans>
\end{lstlisting}

% =============================================================================
% SECTION C.6: TESTING STRATEGIES, TEST FIXTURES & VALIDATION
% =============================================================================
\section{Testing Strategies, Test Fixtures \& Validation}
\label{sec:ergon_testing}

Testing Ergon modules involves isolated unit testing using Apache Camel Test Kit and full container integration testing within Ponos.

\subsection{CamelTestSupport Unit Testing Pattern}

Unit tests extend \texttt{org.apache.camel.test.junit5.CamelTestSupport}, mock downstream completion and error endpoints, inject synthetic \texttt{Pragma} fixtures, and assert exchange transformations and checkpoint states.

\begin{lstlisting}[language=Java, caption={ObservationEnrichmentErgonTest Implementation}]
package net.fhirfactory.harmonia.erga.observation.enrichment;

import ca.uhn.hapi.fhir.context.FhirContext;
import net.fhirfactory.harmonia.calliope.model.ErgonPayload;
import net.fhirfactory.harmonia.calliope.model.Pragma;
import net.fhirfactory.harmonia.calliope.model.PragmaStatus;
import org.apache.camel.EndpointInject;
import org.apache.camel.Produce;
import org.apache.camel.ProducerTemplate;
import org.apache.camel.builder.AdviceWith;
import org.apache.camel.component.mock.MockEndpoint;
import org.apache.camel.test.junit5.CamelTestSupport;
import org.junit.jupiter.api.BeforeEach;
import org.junit.jupiter.api.Test;

import java.util.Date;
import java.util.UUID;

import static org.junit.jupiter.api.Assertions.*;

public class ObservationEnrichmentErgonTest extends CamelTestSupport {

    private ObservationEnrichmentErgon ergon;

    @Produce("direct:activity-ergon-observation-enrichment")
    private ProducerTemplate startEndpoint;

    @EndpointInject("mock:activity-ergon-observation-enrichment-complete")
    private MockEndpoint completeEndpoint;

    @Override
    protected void doPreSetup() throws Exception {
        ergon = new ObservationEnrichmentErgon();
        // Inject field dependencies manually for pure Camel unit testing
        // (In container tests, CDI handles injection)
    }

    @Test
    public void testSuccessfulObservationEnrichment() throws Exception {
        // Arrange
        Pragma syntheticPragma = new Pragma();
        syntheticPragma.setPragmaId("PRAGMA-" + UUID.randomUUID());
        syntheticPragma.setStatus(PragmaStatus.PENDING);
        syntheticPragma.setPatientId("PAT-10045");

        ErgonPayload input = new ErgonPayload();
        input.setPayloadId("PAYLOAD-001");
        input.setPayloadType("HL7_OBX");
        input.setPayload("OBX|1|NM|2823-3^Potassium^LN||5.8|mmol/L|3.5-5.0|H|||F");
        syntheticPragma.addInput(input);

        // Act
        template.sendBody("direct:activity-ergon-observation-enrichment", syntheticPragma);

        // Assert
        assertNotNull(syntheticPragma.getOutputs());
        assertEquals(1, syntheticPragma.getOutputs().size());
        assertEquals("FHIR_R5_OBSERVATION", syntheticPragma.getOutputs().get(0).getPayloadType());
        
        // Assert Checkpoints
        assertFalse(syntheticPragma.getCheckpoints().isEmpty());
        assertEquals("ergon-observation-enrichment", 
                     syntheticPragma.getCheckpoints().get(0).getErgonId());
    }
}
\end{lstlisting}

% =============================================================================
% SECTION C.7: DEPLOYMENT, CONTAINERIZATION & OPERATIONAL RUNBOOK
% =============================================================================
\section{Deployment, Containerization \& Operational Runbook}
\label{sec:ergon_deployment}

\subsection{WildFly \& Ponos Container Architecture}

Ergon modules are bundled into the \texttt{hie-task-processor.war} artifact deployed inside the Ponos WildFly Jakarta EE runtime container. Figure~\ref{fig:ergon_deployment_topo} illustrates the container topology.

\begin{figure}[htbp]
\centering
\begin{lstlisting}[language=XML, caption={Docker Compose Service Definition for Ponos Task Processor}]
  hie-task-processor:
    build:
      context: ./energeia/ponos
      dockerfile: Dockerfile
    container_name: hie-task-processor
    environment:
      - JAVA_OPTS=-Xms1024m -Xmx2048m -XX:+UseG1GC
      - ARTEMIS_HOST=hie-activemq
      - ARTEMIS_PORT=61616
      - INFINISPAN_HOST=hie-infinispan
      - INFINISPAN_PORT=11222
      - DB_URL=jdbc:postgresql://hie-postgres:5432/harmonia_db
      - PONOS_CONCURRENCY_LIMIT=32
    ports:
      - "8080:8080"
      - "9990:9990"
    depends_on:
      - hie-activemq
      - hie-infinispan
      - hie-postgres
    networks:
      - harmonia-mesh
\end{lstlisting}
\caption{Docker Compose Topology for the Ponos Task Processor Container}
\label{fig:ergon_deployment_topo}
\end{figure}

\subsection{Dynamic Live Reload via ponos-cli}

The Ponos container supports zero-downtime dynamic route reloading through the \texttt{ponos-cli} command-line administrative utility. When a new Ergon activity or updated Praxis workflow is packaged into a hot-deployable JAR, operators issue:

\begin{lstlisting}[language=bash, caption={Operating Dynamic Reload with ponos-cli}]
# Check current Ponos worker status and active routes
ponos-cli --status --host localhost --port 8080

# Trigger graceful dynamic route reload
ponos-cli --reload --module ergon-observation-enrichment-1.0.0.jar

# Validate active route registration
ponos-cli --routes | grep "ergon-observation-enrichment"
\end{lstlisting}

\subsection{Concurrency Tuning and Thread Pool Sizing}

Ponos manages Camel concurrency using SEDA and JMS thread pools tuned via deployment parameters:
\begin{itemize}
    \item \textbf{Concurrent Consumers}: Configured on inbound Artemis endpoints (default: 8 consumers, scalable to 32 under peak load).
    \item \textbf{SEDA Queue Capacity}: Internal async handoffs buffer up to 1000 tasks per activity before applying backpressure.
    \item \textbf{Infinispan Cache Write-Behind}: Checkpoints are committed to the Mneme in-memory cache synchronously and flushed to Mnemosyne PostgreSQL asynchronously via write-behind replication queues.
\end{itemize}

\subsection{Dead-Letter Queue (DLQ) Recovery Operations}

When an unrecoverable processing error or transient database failure exhausts configured retry attempts (default: 3 retries with exponential backoff), the message is routed to the Dead-Letter Queue:
\begin{equation}
\texttt{petasos.queue.dlq.ergon-processing}
\end{equation}

Operators inspect and re-drive DLQ tasks using the Hestia Operations CLI:
\begin{lstlisting}[language=bash, caption={DLQ Inspection and Task Redelivery}]
# Inspect dead-lettered Pragma instances
hestia-cli dlq list --queue petasos.queue.dlq.ergon-processing

# View detailed stack trace and error checkpoint
hestia-cli dlq inspect --pragma-id PRAGMA-78a421-99b

# Replay task after remediating downstream dependency
hestia-cli dlq replay --pragma-id PRAGMA-78a421-99b --target-queue petasos.queue.inbound.oru
\end{lstlisting}

\section{Summary \& Best Practices Checklist}
\label{sec:ergon_checklist}

Table~\ref{tab:ergon_checklist} summarizes the mandatory architectural invariants and best practices required when authoring and deploying new Ergon activity modules.

\begin{table}[htbp]
\centering
\small
\begin{tabularx}{\textwidth}{l p{4.5cm} X}
\toprule
\textbf{Phase} & \textbf{Requirement / Invariant} & \textbf{Verification Criteria} \\
\midrule
\textbf{Design} & Inherit from \texttt{ErgonBase} & Verify \texttt{configure()} registers standard input, complete, and error direct endpoints. \\
\textbf{Scoping} & Annotate with \texttt{@Dependent} & Ensure thread-safe RouteBuilder lifecycle per Camel context. \\
\textbf{Validation} & Defensive Inbound Validation & Throw structured \texttt{ErgonPayloadValidationException} on malformed or empty payloads. \\
\textbf{Audit} & Granular \texttt{PragmaCheckpoint} & Record execution timestamps, status, and reason codes in the \texttt{Mneme} grid. \\
\textbf{Security} & Tag Confidentiality Labels & Apply \texttt{FhirSecurityTagManager} classifications to all outgoing FHIR models. \\
\textbf{Testing} & CamelTestSupport Coverage & Author unit tests verifying happy path, error routing, and checkpoint assertions. \\
\textbf{Operations} & Hot Reload Validation & Verify dynamic reload and route registration via \texttt{ponos-cli --reload}. \\
\bottomrule
\end{tabularx}
\caption{Production Readiness Checklist for New Ergon Workflow Activity Modules}
\label{tab:ergon_checklist}
\end{table}
